% ============================================================
% Section 85: 一维晶体能带结构与电子运动
% ============================================================
\section{固体理论：一维晶体能带结构与电子运动}
\label{sec:1d-band-transit-time}

\begin{modelbox}[题目：一维晶体的能带、速度、有效质量与穿越时间]
设一维晶体的电子能带可写成
\[
E(k)=\frac{\hbar^2}{ma^2}\Bigl(\frac{7}{8}-\cos ka+\frac{1}{8}\cos 2ka\Bigr),
\]
其中 $a$ 为晶格常数。
\begin{enumerate}
    \item 求能带的宽度；
    \item 求电子在波矢 $k$ 处的速度；
    \item 求能带底部和顶部电子的有效质量；
    \item 若 $a=2.5$ \AA，当外加 $10^2\,\text{V/m}$ 和 $10^7\,\text{V/m}$ 电场时，分别估计电子自能带底运动到能带顶所需的时间。
\end{enumerate}
\end{modelbox}

\begin{tipbox}[考点快速分析]
\begin{itemize}
    \item \textbf{能带极值}：$dE/dk=0$ 求极值点，比较能量得带宽
    \item \textbf{群速度}：$v=\hbar^{-1}dE/dk$
    \item \textbf{有效质量}：$1/m^*=\hbar^{-2}d^2E/dk^2$，带底正、带顶负
    \item \textbf{$k$ 空间匀速运动}：$t=\Delta k/|dk/dt|=\pi\hbar/(ea\mathcal{E})$
\end{itemize}
\end{tipbox}

\begin{derivationbox}[标准解题过程]
\textbf{Step 1: 能带宽度}

对能量求导：
\begin{align}
\frac{dE}{dk}&=\frac{\hbar^2}{ma^2}\Bigl(a\sin ka-\frac{a}{4}\sin 2ka\Bigr) \notag\\
&=\frac{\hbar^2}{ma}\sin ka\Bigl(1-\frac{1}{2}\cos ka\Bigr).
\end{align}

因 $\cos ka\le1$，括号内恒正，故 $dE/dk=0$ 仅当 $\sin ka=0$。在第一布里渊区 $[-\pi/a,\pi/a]$ 内：

\begin{itemize}
    \item $k=0$：$\cos0=1$，$\cos0=1$
    \begin{equation}
    E(0)=\frac{\hbar^2}{ma^2}\Bigl(\frac{7}{8}-1+\frac{1}{8}\Bigr)=0.\quad\text{（能带底）}
    \end{equation}
    \item $k=\pm\pi/a$：$\cos(\pm\pi)=-1$，$\cos(\pm2\pi)=1$
    \begin{equation}
    E\Bigl(\pm\frac{\pi}{a}\Bigr)=\frac{\hbar^2}{ma^2}\Bigl(\frac{7}{8}+1+\frac{1}{8}\Bigr)=\frac{2\hbar^2}{ma^2}.\quad\text{（能带顶）}
    \end{equation}
\end{itemize}

能带宽度
\begin{equation}
\boxed{\Delta E=E_{\max}-E_{\min}=\frac{2\hbar^2}{ma^2}.}
\end{equation}

\textbf{Step 2: 电子速度}

群速度
\begin{equation}
\boxed{v(k)=\frac{1}{\hbar}\frac{dE}{dk}=\frac{\hbar}{ma}\Bigl(\sin ka-\frac{1}{4}\sin 2ka\Bigr).}
\end{equation}

\textbf{Step 3: 带底和带顶的有效质量}

二阶导数
\begin{equation}
\frac{d^2E}{dk^2}=\frac{\hbar^2}{m}\Bigl(\cos ka-\frac{1}{2}\cos 2ka\Bigr).
\end{equation}

有效质量
\begin{equation}
m^*=\hbar^2\Bigl/\frac{d^2E}{dk^2}=\frac{m}{\cos ka-\frac{1}{2}\cos 2ka}.
\end{equation}

\textbf{带底}（$k=0$）：$\cos0=1$，$\cos0=1$
\begin{equation}
\boxed{m^*(0)=\frac{m}{1-\frac{1}{2}}=2m.}
\end{equation}

\textbf{带顶}（$k=\pi/a$）：$\cos\pi=-1$，$\cos2\pi=1$
\begin{equation}
\boxed{m^*\Bigl(\frac{\pi}{a}\Bigr)=\frac{m}{-1-\frac{1}{2}}=-\frac{2m}{3}.}
\end{equation}

\textbf{Step 4: 电子从带底运动到带顶的时间}

在恒定电场 $\mathcal{E}$ 下，$k$ 空间运动方程
\begin{equation}
\hbar\frac{dk}{dt}=-e\mathcal{E}
\quad\Longrightarrow\quad
\frac{dk}{dt}=-\frac{e\mathcal{E}}{\hbar}.
\end{equation}

从带底 $k=0$ 到带顶 $k=\pi/a$，波矢变化 $\Delta k=\pi/a$。所需时间
\begin{equation}
t=\frac{\Delta k}{|dk/dt|}=\frac{\pi/a}{e\mathcal{E}/\hbar}=\frac{\pi\hbar}{ea\mathcal{E}}.
\end{equation}

代入常数 $\hbar=1.055\times10^{-34}\,\text{J}\cdot\text{s}$，$e=1.6\times10^{-19}\,\text{C}$，$a=2.5\times10^{-10}\,\text{m}$：

\textbf{对 $\mathcal{E}_1=10^2\,\text{V/m}$}：
\begin{equation}
\boxed{t_1=\frac{\pi\times1.055\times10^{-34}}{1.6\times10^{-19}\times2.5\times10^{-10}\times10^2}\approx8.3\times10^{-8}\,\text{s}=83\,\text{ns}.}
\end{equation}

\textbf{对 $\mathcal{E}_2=10^7\,\text{V/m}$}：
\begin{equation}
\boxed{t_2=\frac{\pi\times1.055\times10^{-34}}{1.6\times10^{-19}\times2.5\times10^{-10}\times10^7}\approx8.3\times10^{-13}\,\text{s}=0.83\,\text{ps}.}
\end{equation}
\end{derivationbox}

\begin{formulabox}[答案汇总]
\begin{center}
\begin{tabular}{ll}
\toprule
\textbf{问题} & \textbf{答案} \\
\midrule
(1) 能带宽度 & $\Delta E=2\hbar^2/(ma^2)$ \\
(2) 电子速度 & $v(k)=\dfrac{\hbar}{ma}\bigl(\sin ka-\frac{1}{4}\sin 2ka\bigr)$ \\
(3) 带底有效质量 & $m^*(0)=2m$ \\
(3) 带顶有效质量 & $m^*(\pi/a)=-2m/3$ \\
(4) 运动时间 ($10^2$ V/m) & $t_1\approx8.3\times10^{-8}\,\text{s}$ (83 ns) \\
(4) 运动时间 ($10^7$ V/m) & $t_2\approx8.3\times10^{-13}\,\text{s}$ (0.83 ps) \\
\bottomrule
\end{tabular}
\end{center}
\end{formulabox}

\begin{insightbox}[物理图像]
\begin{itemize}
    \item 能带中的 $\cos2ka$ 项反映\textbf{次近邻跃迁}的贡献。仅保留 $\cos ka$ 时带宽为 $2\hbar^2/(ma^2)$；加入 $\cos2ka$ 后，带底从 $k=0$ 移到 $k=\pm\pi/a$（本题中 $k=0$ 仍为带底，但曲率改变），有效质量从 $m$ 变为 $2m$。
    \item 带底有效质量 $2m$ 大于带顶绝对值 $2m/3$，说明带底比带顶``更平坦''，电子在带底附近更难被加速。
    \item 对于常规电场 $10^2\,\text{V/m}$，穿越时间 83 ns 远大于电子的散射时间（室温下约 $10^{-14}\!\sim\!10^{-12}\,$s），电子在到达带顶之前就被散射，因此实际晶体中观察不到完整的布洛赫振荡。只有在 $10^7\,\text{V/m}$ 的强电场下，穿越时间降至亚皮秒，才可能与散射时间竞争——但此时齐纳隧穿效应也开始变得显著。
\end{itemize}
\end{insightbox}
